% !TeX root = ../main-cn.tex
% UTF-8. Introduction synchronized with paper/论文草稿.md.
\section{引言}
\label{sec:1_introduction}

人体动作生成连接自然语言理解、人体运动建模和虚拟角色控制。
近年来，文本到动作（Text-to-Motion）模型能够根据自然语言生成多样的人体运动，为动画创作、虚拟角色和沉浸式交互提供了便捷的内容生成方式。
从扩散模型 MDM 到采用流匹配的 MotionLab，不同生成架构持续改善文本语义与动作分布的建模能力。
然而，更强的文本生成能力并不直接意味着用户能够控制某一次动作的具体执行方式。

这一困难源于文本描述的欠定性。
例如，“a person walks”指定了行走这一动作类别，却没有唯一确定步速、步频、转向和身体摆动。
即使补充更详细的描述，用户也难以用语言逐帧表达自身运动的连续变化。
因此，如何在保留文本表达高层意图的便利性的同时，引入可直接反映运动实例的控制信号，是文本动作生成走向交互式应用需要解决的问题。

稀疏惯性测量单元（IMU）为这一问题提供了可穿戴的运动信息来源。
头部或手腕设备的加速度及姿态读数能够反映局部身体运动，为生成过程提供文本之外的连续时序约束。
已有 DIP、TransPose、IMUPoser 和 MobilePoser 等方法主要从少量 IMU 估计真实人体姿态，以重建准确度和运动稳定性为主要评价目标。
但局部观测并不能唯一确定全身动作：相近的头部 IMU 读数可能同时对应自然垂臂行走和明显摆臂行走。
文本可以进一步指定未观测身体部位的动作意图，而生成模型提供满足这些条件的全身运动先验。

本文研究 IMU-guided Text-to-Motion Generation，即在文本与稀疏 IMU 的共同条件下生成完整人体动作。
我们关注两种互补关系：当文本描述模糊时，IMU 能否约束具体运动实例；当身体观测稀疏时，文本能否补充未观测部位的动作语义。
该任务保留随机生成能力，允许同一组条件对应多个合理动作，因此既需要衡量与控制目标的接近程度，也需要评价文本语义、生成多样性和动作自然性。
当文本与传感器信号存在冲突时，两种条件如何共同影响输出，同样是需要分析的任务边界。

实现这种互补并非简单地把两个输入拼接起来。
首先，IMU 描述局部的加速度与朝向，而生成模型通常在关节、旋转或潜在动作表示中建模，两者在坐标系、时间分辨率和物理含义上存在差异。
其次，预训练生成模型已具有较强的文本先验，新增控制可能被忽略，也可能过度改变原有生成分布。
最后，约束观测部位与保留未观测部位的语义自由度可能相互竞争，降低局部误差并不必然改善整体动作质量。
这些问题使得跨基座验证、正确与错配条件的对照，以及多随机种子评价成为检验方法有效性的必要环节。

为此，我们提出面向预训练文本到动作模型的轻量条件适配框架 ITM，将传感器信息处理与生成基座的内部计算分离。
框架由三个相互衔接的环节组成：用统一的稀疏 IMU 表示描述可用传感器及其时序读数，通过控制编码器提取运动特征，再由基座适配层将特征接入生成过程。
不同基座共用条件表示与控制学习思路，适配层负责匹配其特征维度、时间分辨率和条件注入位置。
这种模块划分旨在将迁移工作集中于基座适配环节，使数据构建、条件对照和评价流程能够复用。
训练期间冻结预训练生成骨干与文本编码器，仅优化新增的控制模块，以利用已有动作先验并降低参数更新规模。
以 MDM 为例，基座适配层可通过逐层零初始化残差实现控制注入；这一实现体现了统一框架在特定架构上的实例化方式。
我们同时构建 HumanML3D/AMASS 文本、动作与虚拟 IMU 配对数据，为框架训练和跨基座评价提供基础。

框架的目标是在不同架构与不同生成能力的预训练模型上，以可复用的适配流程获得实例级控制，并保留文本语义与动作自然性。
我们以扩散与流匹配生成器为实验载体，从控制有效性、语义互补性和生成质量三个方面考察这一目标。
正确配对、错配与条件关闭的对照用于检验传感器是否提供来自具体运动实例的信息。
固定 IMU 的文本编辑用于考察语言对未观测身体部位的影响，并与固定文本的 IMU 编辑共同分析两类条件的作用。
生成质量评价、损失消融和重复采样进一步用于区分控制收益、随机变化及对原有生成分布的扰动。
这些实验将基座的文本生成能力与新增 IMU 控制能力分别评价，为框架的扩展和适用范围提供依据。

% 严谨实现说明：统一的是条件表示、冻结基座的学习原则和评价流程；当前不是一个控制器 checkpoint 可直接跨所有基座使用。
% MDM 使用逐层残差适配器；MotionLab 使用原生 hint/control token 接口，编码器维度和实现也有差别。
% 不能声称所有基座都采用逐层零初始化注入，或迁移只需修改一个投影层；这些需通过代码接口与迁移成本验证。
% 已完成的跨基座实验为 MDM 与 MotionLab。HY-Motion 尚无 Text+IMU 结果；待完成训练和统一复测后扩展实证主张。
% 严谨结果备忘：MDM 双腕多种子改善稳定，单头改善不显著；MotionLab 改善率尚未达到预设门槛。
% 严谨结果备忘：摆臂、转向响应较明确，速度与步幅不可靠；更强 Text-only 不保证更强实例控制。
% 待补：相同测试 ID、物理时长和重复协议的近期强基座评价，报告各自 Text-only 与 Text+IMU 的质量及控制差值。
% 待补：统一适配接口的实现清单、可训练参数量、显存与迁移改动范围，以支撑“可复用、可扩展”的工程主张。
% 待补：文本编辑下联合比较语义响应与 active-sensor 误差，加入 Text-only 对照；保留失败案例。

本文的贡献概括如下：

\begin{enumerate}
\item 形式化文本与稀疏 IMU 联合条件下的人体动作生成任务，明确实例级运动控制与未观测部位语义补全这两类互补目标。
\item 提出由统一 IMU 表示、时序控制编码和基座适配组成的轻量框架，将预训练生成能力与传感器控制学习分离，并构建配对虚拟 IMU 数据桥。
\item 建立包含正确配对、错配及条件关闭对照的评价方式，在同一基座内固定初始生成噪声，并结合文本编辑与重复采样区分条件响应和有效控制。
\item 通过不同生成架构上的实验，分析新增传感器控制对语义匹配、实例级约束与动作稳定性的影响，为跨基座迁移提供评价依据。
\end{enumerate}

本文面向“语言指定意图、可穿戴设备提供运动细节”的动作生成场景。
通过把预训练生成能力与传感器控制适配分开设计，我们希望明确稀疏 IMU 可以提供哪些增量信息，以及这些信息在不同生成基座中受到怎样的限制。
